\documentclass[12pt]{article}
\usepackage[utf8]{inputenc}
\usepackage[T1]{fontenc}
\usepackage{hyperref}
\title{Lower Tanana Language Resources for K-12 Curriculum Development}
\author{Compiled by Claude Sonnet 5 (claude-sonnet-5) for Andr\'as Kornai}
\date{}
\begin{document}
\maketitle

\section{Introduction}
Lower Tanana (ISO 639-3 \texttt{taa}, Glottocode lowe1425) is an Athabaskan language of the middle Tanana River valley in Alaska. It is spoken by about 15 elderly people and is not known to be taught in schools. The search found mainly dictionaries and lexical documentation. It found no school primers, textbooks or curriculum materials.

\section{Primary Educational Resources}
The search found no primers or textbooks for classroom use. The pocket dictionary below is the closest item to a learner's resource.
\begin{itemize}
\item Benhti Kokht'ana Kenaga': Minto Lower Tanana Athabascan Pocket Dictionary (Alaska Native Language Center, 2009; the online version is a draft revision): \url{https://knowledgebase.arts.ubc.ca/benhti-kokhtana-kenaga-minto-lower-tanana-athabascan-pocket-dictionary}
\end{itemize}

\section{Linguistic and Reference Materials}
\begin{itemize}
\item Lower Tanana Athabaskan Dictionary (James Kari, 1994 preliminary draft, a bilingual Lower Tanana--English dictionary with an English index), described in the UBC Relational Lexicography knowledgebase: \url{https://knowledgebase.arts.ubc.ca/lower-tanana-athabaskan-dictionary/}
\item The same dictionary, UBC wiki entry: \url{https://wiki.ubc.ca/Documentation:RelLex/Lower_Tanana_Athabaskan_Dictionary}
\item Minto-Nenana Athabaskan Noun Dictionary (Michael E. Krauss, 1974), a topical English-to-Lower Tanana noun dictionary: \url{https://knowledgebase.arts.ubc.ca/minto-nenana-athabaskan-noun-dictionary}
\item Lower Tanana Dene Dictionary (918 pages). The search result describes it as a consolidation of nearly all primary data from 1903 to 2023: \url{https://www.ubcpress.ca/lower-tanana-dene-dictionary}
\end{itemize}

\section{Religious and Community Resources}
The search found none.

\section{Digital and Audio Resources}
The dictionaries above are described as open access through the Alaska Native Language Archive. I did not verify the archive item pages themselves. The search found no standalone audio or software.

\section{Curriculum Development Context}
Very little classroom material exists, and the speaker community is very small. A curriculum would have to be built from the dictionaries, which would need pedagogical adaptation and work with elders and the language center. Current availability of the printed items should be checked.
\end{document}